% Topic 1.1.01 — Gaussian Elimination.

\section{Gaussian Elimination}
\label{sec:1.1.01-systems-gauss}

\begin{ticket}
Systems of linear equations. Rectangular matrices. Reducing matrices and systems of linear equations to row echelon form. Gaussian elimination.
\end{ticket}

\subsection*{Theory}

We work over a fixed field~$\F$ (think $\F = \R$ throughout). Letters
$i, j, k$ denote integer indices and $m, n$ are positive integers.

\begin{definition}
\label{def:gauss-system}
A \emph{system of $m$ linear equations in $n$ unknowns} over~$\F$ is a list
\[
  \begin{aligned}
    a_{11} x_1 + a_{12} x_2 + \cdots + a_{1n} x_n &= b_1, \\
    &\;\;\vdots \\
    a_{m1} x_1 + a_{m2} x_2 + \cdots + a_{mn} x_n &= b_m,
  \end{aligned}
\]
where the coefficients $a_{ij}$ and right-hand sides $b_i$ lie in~$\F$.
Collecting them gives the \emph{coefficient matrix}
$\mat{A} = (a_{ij}) \in \F^{m \times n}$, the column $\vect{x} =
(x_1, \dots, x_n)\T$ of unknowns, and the column $\vect{b} = (b_1, \dots, b_m)\T$
of right-hand sides; the system is then $\mat{A}\vect{x} = \vect{b}$. The
\emph{augmented matrix} $[\mat{A} \mid \vect{b}] \in \F^{m \times (n+1)}$
records both at once.

A \emph{solution} is a vector $\vect{x} \in \F^n$ that satisfies every
equation. The system is \emph{consistent} if its solution set is nonempty.
\end{definition}

\begin{definition}
\label{def:gauss-ref}
A matrix is in \emph{row echelon form} (REF) when:
\begin{enumerate}[label=(\alph*)]
  \item every all-zero row sits below every nonzero row;
  \item in any two consecutive nonzero rows, the leading nonzero entry
        (the \emph{pivot}) of the lower row is strictly to the right of
        the pivot of the upper one.
\end{enumerate}
It is in \emph{reduced row echelon form} (RREF) if additionally:
\begin{enumerate}[label=(\alph*),resume]
  \item every pivot equals~$1$;
  \item each pivot is the only nonzero entry in its column.
\end{enumerate}
The columns containing pivots are \emph{pivot columns}; the others are
\emph{free columns}. Variables corresponding to pivot columns are
\emph{pivot variables}; those corresponding to free columns are
\emph{free variables}.
\end{definition}

\begin{definition}
\label{def:gauss-ero}
The three \emph{elementary row operations} (EROs) on a matrix are:
\begin{enumerate}[label=(\Roman*)]
  \item swap rows $i$ and~$j$, written $R_i \leftrightarrow R_j$;
  \item scale row $i$ by a nonzero $c \in \F$, written $R_i \to c\,R_i$;
  \item add $c$ times row $j$ to row $i$ (with $i \neq j$), written
        $R_i \to R_i + c\,R_j$.
\end{enumerate}
Each ERO is realized as left multiplication by an \emph{elementary matrix}:
the identity $I_m$ with one row modified by the same operation. Every
elementary matrix is invertible, and its inverse is again elementary
(swap is self-inverse; scale by $c$ inverts to scale by $c^{-1}$;
$R_i \to R_i + c\,R_j$ inverts to $R_i \to R_i - c\,R_j$).
\end{definition}

\begin{theorem}[Gauss]
\label{thm:gauss-rref}
Every matrix $\mat{A} \in \F^{m \times n}$ can be transformed into a
matrix in RREF by a finite sequence of EROs. Moreover, the RREF
of~$\mat{A}$ is uniquely determined.
\end{theorem}
\proofof{gauss-rref}

\begin{theorem}[Invariance of solutions under EROs]
\label{thm:gauss-invariance}
If $[\mat{A}' \mid \vect{b}']$ is obtained from $[\mat{A} \mid \vect{b}]$
by a finite sequence of EROs, then $\mat{A}\vect{x} = \vect{b}$ and
$\mat{A}'\vect{x} = \vect{b}'$ have the same solution set.
\end{theorem}
\proofof{gauss-invariance}

\begin{corollary}[Solving via Gaussian elimination]
\label{cor:gauss-solving}
To solve $\mat{A}\vect{x} = \vect{b}$:
\begin{enumerate}
  \item form the augmented matrix $[\mat{A} \mid \vect{b}]$;
  \item reduce it to REF (or RREF) by EROs; call the result
        $[\mat{U} \mid \vect{c}]$;
  \item the system is \emph{inconsistent} iff some row of $[\mat{U} \mid \vect{c}]$
        has the form $(0, \dots, 0 \mid c)$ with $c \neq 0$;
  \item otherwise, treat each free variable as a parameter and recover
        the pivot variables one by one, working from the bottommost
        pivot row upward (\emph{back-substitution}).
\end{enumerate}
\end{corollary}
\proofof{gauss-solving}

\begin{remark}
\label{rem:gauss-complexity}
Reducing an $m \times n$ matrix to REF by the algorithm in the proof of
\cref{thm:gauss-rref} costs $O(mn \cdot \min(m, n))$ field operations;
in particular, $O(n^3)$ for square systems. The uniqueness clause of
\cref{thm:gauss-rref} is best handled with the language of rank and
column spaces and is revisited in \cref{sec:1.1.02-rank}. For a
fuller treatment of the standard linear-algebra background, see
\cite{Beklemishev2025}.
\end{remark}

\subsection*{Proofs}

\begin{proof}[\Cref{thm:gauss-invariance}]
\proofanchor{gauss-invariance}
Write $\vect{r}_i \in \F^n$ for the $i$-th row of~$\mat{A}$, so the
$i$-th equation of the system reads $\vect{r}_i \cdot \vect{x} = b_i$.
We check that each ERO sends solutions to solutions; reversibility of
EROs (\cref{def:gauss-ero}) then gives the converse.

\emph{Type~(I), swap.} Permuting the equations leaves the conjunction
unchanged.

\emph{Type~(II), scale.} If $\vect{r}_i \cdot \vect{x} = b_i$ and
$c \neq 0$, multiplying both sides by $c$ gives
$(c\,\vect{r}_i) \cdot \vect{x} = c\,b_i$. The reverse implication
follows by multiplying by $c^{-1}$.

\emph{Type~(III), replace.} Suppose $\vect{r}_i \cdot \vect{x} = b_i$
and $\vect{r}_j \cdot \vect{x} = b_j$. Adding $c$ times the second
equation to the first gives
$(\vect{r}_i + c\,\vect{r}_j) \cdot \vect{x} = b_i + c\,b_j$. The
reverse implication uses $R_i \to R_i - c\,R_j$.

In all three cases the new equation is satisfied by exactly the same
$\vect{x}$. Composing finitely many EROs preserves the solution set.
\end{proof}

\begin{proof}[\Cref{thm:gauss-rref}, existence]
\proofanchor{gauss-rref}
We give the algorithm and verify termination.

\emph{Phase~1 (forward elimination, REF).} If $\mat{A}$ is the zero
matrix, it is already in REF. Otherwise let $j_1$ be the leftmost
column containing a nonzero entry, and pick a row $i_1$ with
$a_{i_1, j_1} \neq 0$. Swap rows $1$ and $i_1$ (operation~I) to bring
this entry into position $(1, j_1)$ — call it the first \emph{pivot}.
For each row $i > 1$, apply
$R_i \to R_i - (a_{i, j_1}/a_{1, j_1})\,R_1$ (operation~III) to zero out
the entry in column~$j_1$. The submatrix formed by rows $2, \dots, m$
and columns $j_1 + 1, \dots, n$ is now strictly smaller (one row fewer
and at least one column fewer), so the same procedure applied
recursively terminates.

\emph{Phase~2 (back substitution, RREF).} Starting from REF, scale each
pivot row by the inverse of its pivot (operation~II) so every pivot
becomes~$1$. Then process the pivot rows from bottom to top: for the
$k$-th pivot, in column $j_k$, run $R_i \to R_i - a_{i, j_k}\,R_k$
(operation~III) for each $i < k$ to clear all entries above the pivot.
Each step preserves the REF shape and strictly reduces the count of
nonzero entries above pivots, so the process terminates after at most
$\binom{r}{2}$ steps, where $r$ is the number of pivots.

The result is in RREF, established by a finite sequence of EROs.

Uniqueness of the RREF is deferred to \cref{sec:1.1.02-rank}, where
it follows from the invariance of the row space.
\end{proof}

\begin{proof}[\Cref{cor:gauss-solving}]
\proofanchor{gauss-solving}
By \cref{thm:gauss-invariance}, the reduced augmented matrix
$[\mat{U} \mid \vect{c}]$ has exactly the same solution set as the
original system.

A row of the form $(0, \dots, 0 \mid c)$ with $c \neq 0$ encodes the
equation $0 = c$, which has no solutions; if any such row is present
the system is inconsistent. Conversely, if no such row is present, then
every nonzero row of $[\mat{U} \mid \vect{c}]$ has its pivot in some
column $j_k \le n$ of the coefficient block, and the corresponding
equation determines the pivot variable $x_{j_k}$ as an affine function
of the variables in columns $> j_k$ (i.e.\ free variables further to
the right and pivot variables not yet recovered). Resolving from the
last pivot row up assigns $x_{j_k}$ for each~$k$ and leaves free
variables as parameters; the resulting parametric family is the entire
solution set.
\end{proof}

\subsection*{Examples}

\begin{example}[Unique solution]
\label{ex:gauss-unique}
Solve
\[
  \begin{aligned}
    x + y + z &= 6, \\
    2x + y - z &= 1, \\
    3x - y + 2z &= 7.
  \end{aligned}
\]
Apply Gaussian elimination to the augmented matrix:
\[
  \left[\begin{array}{rrr|r}
    1 &  1 &  1 & 6 \\
    2 &  1 & -1 & 1 \\
    3 & -1 &  2 & 7
  \end{array}\right]
  \xrightarrow{\substack{R_2 \to R_2 - 2R_1 \\ R_3 \to R_3 - 3R_1}}
  \left[\begin{array}{rrr|r}
    1 &  1 &  1 &   6 \\
    0 & -1 & -3 & -11 \\
    0 & -4 & -1 & -11
  \end{array}\right]
  \xrightarrow{R_3 \to R_3 - 4R_2}
  \left[\begin{array}{rrr|r}
    1 &  1 &  1 &   6 \\
    0 & -1 & -3 & -11 \\
    0 &  0 & 11 &  33
  \end{array}\right].
\]
The matrix is in REF; back-substitution gives $z = 3$, then
$-y - 9 = -11$, so $y = 2$, and finally $x + 2 + 3 = 6$, so $x = 1$.
The unique solution is $\vect{x} = (1, 2, 3)\T$.
\end{example}

\begin{example}[Two-parameter solution family]
\label{ex:gauss-parametric}
Solve
\[
  \begin{aligned}
    x + 2y + z + w &= 4, \\
    2x + 4y + z - w &= 5.
  \end{aligned}
\]
One ERO suffices to reach REF:
\[
  \left[\begin{array}{rrrr|r}
    1 & 2 &  1 &  1 & 4 \\
    2 & 4 &  1 & -1 & 5
  \end{array}\right]
  \xrightarrow{R_2 \to R_2 - 2R_1}
  \left[\begin{array}{rrrr|r}
    1 & 2 &  1 &  1 & 4 \\
    0 & 0 & -1 & -3 & -3
  \end{array}\right].
\]
Pivots are in columns~$1$ and~$3$, so $x$ and $z$ are pivot variables
and $y, w$ are free. From the second row $z = 3 - 3w$, and from the
first $x = 4 - 2y - z - w = 1 - 2y + 2w$. The solution set is
\[
  \begin{pmatrix} x \\ y \\ z \\ w \end{pmatrix}
  =
  \begin{pmatrix} 1 \\ 0 \\ 3 \\ 0 \end{pmatrix}
  + y \begin{pmatrix} -2 \\ 1 \\ 0 \\ 0 \end{pmatrix}
  + w \begin{pmatrix} 2 \\ 0 \\ -3 \\ 1 \end{pmatrix},
  \qquad y, w \in \F.
\]
A particular solution plus a two-dimensional family parametrized by the
two free variables.
\end{example}

\begin{example}[Inconsistent system]
\label{ex:gauss-inconsistent}
The system
\[
  \begin{aligned}
    x + 2y &= 3, \\
    2x + 4y &= 7
  \end{aligned}
\]
reduces by $R_2 \to R_2 - 2R_1$ to
\[
  \left[\begin{array}{rr|r}
    1 & 2 & 3 \\
    0 & 0 & 1
  \end{array}\right],
\]
whose second row encodes $0 = 1$. By \cref{cor:gauss-solving} the
system has no solutions: the second equation contradicts the first.
\end{example}

\begin{problem}
\label{prob:gauss-parameter}
For which values of $k \in \R$ is the system
\[
  \begin{aligned}
    x + y + z &= 1, \\
    x + 2y + 4z &= k, \\
    x + 4y + 10z &= k^2
  \end{aligned}
\]
(a)~inconsistent, (b)~uniquely solvable, (c)~consistent with infinitely
many solutions? Describe the solution set in case~(c).
\end{problem}

\begin{proof}[Solution]
Eliminate using the first row:
\[
  \left[\begin{array}{rrr|c}
    1 & 1 &  1 & 1 \\
    1 & 2 &  4 & k \\
    1 & 4 & 10 & k^2
  \end{array}\right]
  \xrightarrow{\substack{R_2 \to R_2 - R_1 \\ R_3 \to R_3 - R_1}}
  \left[\begin{array}{rrr|c}
    1 & 1 & 1 & 1     \\
    0 & 1 & 3 & k - 1 \\
    0 & 3 & 9 & k^2 - 1
  \end{array}\right]
  \xrightarrow{R_3 \to R_3 - 3R_2}
  \left[\begin{array}{rrr|c}
    1 & 1 & 1 & 1     \\
    0 & 1 & 3 & k - 1 \\
    0 & 0 & 0 & (k - 1)(k - 2)
  \end{array}\right],
\]
since $k^2 - 1 - 3(k - 1) = k^2 - 3k + 2 = (k - 1)(k - 2)$.

The third row is $0 = (k - 1)(k - 2)$. Hence:
\begin{itemize}
  \item if $k \notin \{1, 2\}$, the right-hand side is nonzero and the
        system is \emph{inconsistent};
  \item if $k \in \{1, 2\}$, the third row vanishes; the matrix has
        only two pivots in three columns, so~(b) never occurs and
        case~(c) applies, with $z$ as the single free variable.
\end{itemize}
For $k = 1$: $y + 3z = 0$, so $y = -3z$, and $x + y + z = 1$ gives
$x = 1 + 2z$; the solution set is
$(1, 0, 0)\T + z\,(2, -3, 1)\T$ for $z \in \R$.

For $k = 2$: $y + 3z = 1$, so $y = 1 - 3z$, and $x = 1 - y - z = 2z$;
the solution set is $(0, 1, 0)\T + z\,(2, -3, 1)\T$ for $z \in \R$.
\end{proof}
